% !TeX program = xelatex
% !TeX encoding = UTF-8
\documentclass[UTF8,zihao=-4,fontset=windows]{ctexart}

% ---------- 页面与字体（像 Word 的“页面布局 / 字体”） ----------
\usepackage[a4paper,top=2.54cm,bottom=2.54cm,left=3.00cm,right=2.50cm]{geometry}
\usepackage{fontspec}
\setmainfont{Times New Roman}
\setCJKmainfont[BoldFont=SimHei]{SimSun}
\setCJKsansfont{Microsoft YaHei}

% ---------- 段落（像 Word 的“段落”对话框） ----------
\usepackage{setspace}
\setstretch{1.5}                 % 1.5 倍行距
\setlength{\parindent}{2em}      % 首行缩进 2 个汉字
\setlength{\parskip}{0pt}        % 段后间距
\usepackage{ragged2e}
\AtBeginDocument{\justifying}   % 两端对齐

% ---------- 标题、页眉页脚、图表与数学 ----------
\ctexset{
  section={format=\heiti\zihao{3},beforeskip=18pt,afterskip=10pt},
  subsection={format=\heiti\zihao{4},beforeskip=14pt,afterskip=7pt},
  subsubsection={format=\heiti\zihao{-4},beforeskip=10pt,afterskip=5pt}
}
\usepackage{amsmath,amssymb}
\usepackage{graphicx}
\usepackage{booktabs}
\usepackage{caption}
\usepackage{subcaption}
\usepackage{float}
\captionsetup{font=small,labelsep=quad}
\usepackage{enumitem}
\setlist{nosep,leftmargin=2em}
\usepackage[hidelinks]{hyperref}
\usepackage{fancyhdr}
\pagestyle{fancy}
\fancyhf{}
\fancyhead[C]{\zihao{5} LaTeX 论文写作示例}
\fancyfoot[C]{\thepage}
\setlength{\headheight}{14pt}

% Figure 演示用占位框：换成 \includegraphics 即可使用真实图片。
\newcommand{\figureplaceholder}[3][4cm]{%
  \fbox{\parbox[c][#1][c]{\dimexpr#2-2\fboxsep-2\fboxrule\relax}{%
    \centering\sffamily\zihao{5}#3}}
}

% ---------- 参考文献 ----------
\usepackage[backend=biber,style=numeric,sorting=none]{biblatex}
\addbibresource{references.bib}

\title{LaTeX 中文论文写作示例}
\author{作者姓名}
\date{\today}

\begin{document}
\maketitle

\begin{abstract}
这是一份可直接编译的中文论文起始模板。正文默认小四号、1.5 倍行距、首行缩进两个汉字，并展示标题、对齐、公式、表格和参考文献。将这里替换为你的摘要。
\end{abstract}

\section{像 Word 一样控制排版}

正文的全局字号由文档类选项 \verb|zihao=-4| 控制，即小四号。局部字号和精确磅值可以分别写成：
\begin{flushleft}
  \verb|{\zihao{5} 五号文字}|\\
  \verb|{\fontsize{10.5pt}{16pt}\selectfont 文字}|
\end{flushleft}
精确字号命令中的第二个数字是基线间距。

段落间距、首行缩进和行距集中写在导言区，通常不应逐段手调。需要取消缩进时使用
\verb|\noindent|。居中、左对齐和右对齐分别使用 \verb|center|、\verb|flushleft|、
\verb|flushright| 环境。

\begin{center}
  这是一行居中的文字。
\end{center}

\subsection{公式与编号}

行内公式例如 $E=mc^2$。独立且自动编号的公式如下：
\begin{equation}
  f(x)=\int_{-\infty}^{\infty}\hat f(\xi)e^{2\pi i x\xi}\,\mathrm{d}\xi.
  \label{eq:fourier}
\end{equation}
上式可以被交叉引用；在 LaTeX 中使用 \verb|\label| 标记，再用 \verb|\eqref| 引用即可。

\subsection{表格}

\begin{table}[htbp]
  \centering
  \caption{常见中文字号与磅值近似对应}
  \label{tab:font-size}
  \begin{tabular}{ccc}
    \toprule
    中文字号 & LaTeX 命令 & 近似磅值 \\
    \midrule
    三号 & \verb|\zihao{3}| & 16 pt \\
    小四 & \verb|\zihao{-4}| & 12 pt \\
    五号 & \verb|\zihao{5}| & 10.5 pt \\
    小五 & \verb|\zihao{-5}| & 9 pt \\
    \bottomrule
  \end{tabular}
\end{table}

\section{Figure 布局演示}

\subsection{单图：按栏宽缩放}

\begin{figure}[htbp]
  \centering
  % 真实图片：\includegraphics[width=0.82\linewidth]{figures/result.pdf}
  \figureplaceholder[4.5cm]{0.82\linewidth}{单图占位\\[4pt]
    \texttt{width=0.82\textbackslash linewidth}}
  \caption{单图居中，宽度为当前栏宽的 82\%}
  \label{fig:single}
\end{figure}

\subsection{并排双图：各自标题与总标题}

\begin{figure}[htbp]
  \centering
  \begin{subfigure}[t]{0.48\linewidth}
    \centering
    % \includegraphics[width=\linewidth]{figures/method-a.pdf}
    \figureplaceholder[3.6cm]{\linewidth}{子图 A}
    \caption{方法 A}
    \label{fig:pair-a}
  \end{subfigure}\hfill
  \begin{subfigure}[t]{0.48\linewidth}
    \centering
    % \includegraphics[width=\linewidth]{figures/method-b.pdf}
    \figureplaceholder[3.6cm]{\linewidth}{子图 B}
    \caption{方法 B}
    \label{fig:pair-b}
  \end{subfigure}
  \caption{两个等宽子图；\texttt{\textbackslash hfill} 自动分配中间空隙}
  \label{fig:pair}
\end{figure}

引用整个图用“图~\ref{fig:pair}”，引用其中一幅用“图~\ref{fig:pair-a}”。

\subsection{2×2 子图网格}

\begin{figure}[H]
  \centering
  \begin{subfigure}[t]{0.48\linewidth}
    \centering
    \figureplaceholder[2.5cm]{\linewidth}{子图 A}
    \caption{数据集 A}
  \end{subfigure}\hfill
  \begin{subfigure}[t]{0.48\linewidth}
    \centering
    \figureplaceholder[2.5cm]{\linewidth}{子图 B}
    \caption{数据集 B}
  \end{subfigure}

  \medskip
  \begin{subfigure}[t]{0.48\linewidth}
    \centering
    \figureplaceholder[2.5cm]{\linewidth}{子图 C}
    \caption{数据集 C}
  \end{subfigure}\hfill
  \begin{subfigure}[t]{0.48\linewidth}
    \centering
    \figureplaceholder[2.5cm]{\linewidth}{子图 D}
    \caption{数据集 D}
  \end{subfigure}
  \caption{四个子图组成的 2×2 网格}
  \label{fig:grid}
\end{figure}

浮动参数 \verb|[htbp]| 表示依次尝试当前位置、页顶、页底和单独浮动页。双栏论文中，
普通 \verb|figure| 只占一栏；把环境改成 \verb|figure*| 可跨两栏，通常出现在页顶。

\section{引用与投稿格式}

LaTeX 的优势是将内容与版式分离。投稿时通常替换 \verb|\documentclass|，并按期刊要求加载其
\verb|.cls| 或官方模板，不必逐段重排。参考文献可以像这样引用经典文献~\cite{knuth1984}。

\printbibliography
\end{document}
